\documentclass[lettersize,journal]{IEEEtran}
\usepackage{amsmath,amsfonts}
\usepackage{algorithm}
\usepackage{array}
\usepackage[caption=false,font=normalsize,labelfont=sf,textfont=sf]{subfig}
\usepackage{textcomp}
\usepackage{stfloats}
\usepackage{url}
\usepackage{verbatim}
\usepackage{graphicx}
\usepackage{cite}
\usepackage{algpseudocode}
\usepackage{amsthm}


\newtheorem{definition}{Definition}
\newtheorem{proposition}{Proposition}
\newtheorem{assumption}{Assumption}
\newtheorem{remark}{Remark}
\newcommand{\Game}{\mathfrak{G}}
\newcommand{\StateSpace}{\mathcal{S}}
\newcommand{\ObsSpace}{\mathcal{O}}
\newcommand{\ActionSpace}{\mathcal{A}}
\newcommand{\TransitionKernel}{\mathbb{P}}
\newcommand{\Payoff}{u}
\newcommand{\DiscFactor}{\gamma}
\newcommand{\FeasGraph}{\mathcal{G}^{\mathrm{feas}}}
\newcommand{\CollabGraph}{\mathcal{G}^{\mathrm{coll}}}
\newcommand{\LinkMap}{\Psi}

\algrenewcommand\algorithmiccomment[1]{\hfill\(\triangleright\) \textit{#1}}
\algrenewcommand\algorithmicrequire{\textbf{Input:}}
\algrenewcommand\algorithmicensure{\textbf{Output:}}


\hyphenation{op-tical net-works semi-conduc-tor IEEE-Xplore}
% updated with editorial comments 8/9/2021

\begin{document}

\title{Dynamic Network Formation for Robust Decentralized Personalized Learning over Hybrid 5G C-V2X Vehicular Networks}

\author{IEEE Publication Technology,~\IEEEmembership{Staff,~IEEE,}
	% <-this % stops a space
	\thanks{This paper was produced by the IEEE Publication Technology Group. They are in Piscataway, NJ.}% <-this % stops a space
	\thanks{Manuscript received April 19, 2021; revised August 16, 2021.}}

% The paper headers
\markboth{Journal of \LaTeX\ Class Files,~Vol.~14, No.~8, August~2021}%
{Shell \MakeLowercase{\textit{et al.}}: A Sample Article Using IEEEtran.cls for IEEE Journals}

\IEEEpubid{0000--0000/00\$00.00~\copyright~2021 IEEE}
% Remember, if you use this you must call \IEEEpubidadjcol in the second
% column for its text to clear the IEEEpubid mark.

\maketitle

\begin{abstract}
	Connected vehicles increasingly rely on collaborative intelligence for services such as cooperative perception, risk prediction, and personalized driver assistance. Realizing such intelligence requires learning from distributed, privacy-sensitive, and strongly non-IID data under stringent latency, energy, and bandwidth constraints. Existing federated and edge-assisted methods remain poorly suited to this setting: they rely on centralized orchestration, assume a single communication interface, and offer limited resilience to unreliable peers. We study robust decentralized personalized learning over hybrid 5G C-V2X networks, whose PC5 sidelink and Uu interfaces offer complementary range, latency, and cost profiles for resource-aware collaboration. We formulate collaboration as a dynamic directed network formation game in which each vehicle chooses its collaborators, communication mode, and aggregation weights under local resource and trust constraints. Based on this formulation, we propose \textsc{Dante}, a fully decentralized framework in which every vehicle learns a link-formation policy through an attention-guided actor--critic, aggregates received models under fused attention--trust--robustness weights, and maintains online local trust estimates for its peers. We analyze convergence and per-round complexity, and SUMO-driven experiments over real urban and highway maps show that \textsc{Dante} delivers substantial resource savings and more efficient collaboration while preserving personalized accuracy and Byzantine robustness over representative baselines.
\end{abstract}

\begin{IEEEkeywords}
	Decentralized federated learning, personalization, vehicular networks, 5G C-V2X, network formation, Byzantine robustness.
\end{IEEEkeywords}


\section{Introduction}
\label{sec:introduction}

Modern connected vehicles are increasingly equipped with high resolution cameras, LiDAR, radar, telemetry modules, and driver or cabin monitoring systems that continuously observe both the external traffic scene and the internal vehicle context~\cite{rajashekara2024data,li2020smart}. Consequently, vehicles generate massive volumes of multimodal data, with recent surveys reporting data rates that can reach the terabyte per hour scale in advanced autonomous driving settings~\cite{rajashekara2024data,li2020smart}. Exploiting such data enables a broad class of learning tasks, including collision risk estimation, trajectory forecasting, cooperative perception, and personalized driver assistance, all of which require models that are both context aware and responsive to local driving conditions~\cite{yu2022personalized}.

Traditionally, such models have been trained by uploading vehicle data to cloud servers, where a global model is learned over the aggregated corpus. This paradigm is impractical in current vehicular environments as transferring massive amounts of data is prohibitively expensive both in terms of bandwidth and energy. Furthermore, it conflicts with strict latency constraints and automotive privacy regulations.
Federated learning (FL) partially addresses this tension by enabling collaborative training through model exchange rather than raw data sharing~\cite{mcmahan2017communication,lim2020survey}. Classical FL, however, remains fundamentally server centric. Its dependence on a central aggregator introduces a communication bottleneck and a single point of failure, while the learned objective is still biased toward a common model that may not adequately serve vehicles operating under distinct traffic regimes, road geometries, sensing conditions, and driver preferences.


To reduce dependence on remote cloud coordination, several vehicular learning systems have shifted aggregation to roadside units, base stations, or edge servers near the road network~\cite{wang2023mobilityaware,pervej2023resource,liao2025pfedvs}. Although this architecture lowers communication latency, it does not remove central orchestration. Vehicles must still maintain infrastructure connectivity, contend for shared uplink resources, and often contribute to regional or global models whose inductive bias remains insufficiently personalized. Mobility further complicates this design because short contact times, frequent handovers, and intermittent coverage can degrade both training efficiency and aggregation quality. More importantly, infrastructure centric coordination underuses a natural source of learning efficiency in vehicular environments, namely direct collaboration among vehicles that experience similar local contexts and can often exchange updates more cheaply than through a base station.

A more compelling substrate for collaboration is provided by hybrid 5G C-V2X networks, standardized by 3GPP~\cite{3gpp2021architecture}. In 3GPP compliant V2X systems, vehicles can communicate directly over the PC5 sidelink, operating in the 5.9\,GHz ITS band, to exchange messages directly with nearby vehicles at low latency and minimal infrastructure cost~\cite{molina2017ltev}. When longer-range communication is required, the Uu interface is used through the 5G cellular network at higher cost but broader reach~\cite{garcia2021tutorial}, thereby exposing two communication modes with markedly different latency, range, and resource profiles. This dual-interface capability creates a natural substrate for decentralized vehicle-to-vehicle (V2V) learning: nearby vehicles can collaborate efficiently and frequently share similar environment context, yet the most important collaborators may not always be the nearest ones. The PC5 sidelink can therefore support low latency exchange among adjacent vehicles, whereas the Uu path can selectively reach statistically useful but geographically distant collaborators through the cellular infrastructure. At the same time, local datasets are non-IID across vehicles, as each reflects a distinct driving environment. Effective collaboration must consequently be selective, personalized, and context aware.

Recent works in the collaborative learning literature~\cite{kharrat2025dpfl,liu2024dfedpgp,ye2023pfedgraph,chen2022graph,zhang2021fedfomo,ma2021pfedatt,zantedeschi2020fully} show that personalization and topology structure are key elements in achieving per client satisfactory models. Many approaches explored in the literature demonstrate that carefully selected collaborators can substantially improve vehicle specific performance under heterogeneous data and constrained communication. Yet the assumptions behind these approaches remain largely misaligned with vehicular networks. Most operate over comparatively stable client populations, infer collaboration in a logically centralized manner, or assume a single communication interface. Existing vehicular decentralized schemes exhibit the complementary limitation: they improve training efficiency under mobility and wireless constraints, but they do not jointly address personalization, hybrid link heterogeneity, or Byzantine robustness within a unified topology control problem~\cite{chen2025mobility,meng2025resource,yuan2023p2p}. In this setting, the collaboration topology is not a networking detail. It is a learning variable because the value of a collaborator depends jointly on vehicle similarity, environment context, link cost, and vehicle trustworthiness.

Vehicular neighborhoods are highly dynamic, as vehicles move rapidly across space and maintain only short-lived communication opportunities. Each vehicle must therefore learn under strict bandwidth, energy, and time constraints. Additionally, vehicular learning takes place in an open environment characterized by strong data heterogeneity and potentially unreliable participants. Personalization and Byzantine resilience are therefore essential requirements in V2X environments. Indiscriminate collaboration with many peers can waste scarce resources, increase computation overhead, and even degrade learning performance when selected neighbors are statistically irrelevant or malicious. These challenges motivate the central question of this paper: \emph{how should a vehicle in a hybrid 5G C-V2X network decide, at each stage of training, which neighbors to collaborate with and over which interface, so as to maximize personalized learning benefit under mobility, resource, and trust constraints without relying on centralized coordination?}


We answer this question by proposing \textsc{DANTE} (\textbf{D}ual-interface \textbf{A}ttention-guided \textbf{N}etwork-formation for \textbf{T}rustworthy vehicular \textbf{E}dge learning), a robust decentralized personalized learning framework built on dynamic network formation over a hybrid V2X collaboration graph. Unlike conventional vehicular FL, which treats communication as a mere transport layer, DANTE jointly models statistical relevance, radio reachability, and trust in topology control. Each vehicle maintains a personalized model and acts as a strategic agent that selects collaborators and assigns aggregation weights. The utility of each candidate link captures expected personalized loss reduction, communication and computation cost, and trust-aware robustness. Short-range exchanges are carried over the PC5 sidelink, while the Uu interface is used selectively to reach high-value collaborators beyond direct radio range. To adapt to rapidly changing neighborhoods, each vehicle learns its link formation policy through a local attention-guided PPO mechanism that encodes candidate neighbors, link quality, budget state, and past collaboration quality, and activates only links expected to improve personalization per unit cost. Training is fully decentralized: no central critic, shared parameters, or global state is required, and each agent learns solely from local observations and local reward. Robust aggregation and online trust estimation are integrated into the same loop, allowing vehicles to mitigate suspicious updates and revert to local training when collaboration becomes unreliable.

The main contributions of this paper are summarized as follows.
\begin{enumerate}
\item We formulate robust decentralized personalized learning over hybrid 5G C-V2X as a directed network formation game in which vehicles are strategic players, link formation decisions define the action space, and payoffs integrate personalized loss improvement, interface aware resource cost, and trust weighted collaboration quality.
\item Each vehicle independently learns a \textsc{Dante} policy for neighbor selection and resource allocation from local observations and local reward, without any central critic, shared parameters, or global state, consistent with the operational constraints of V2X networks.
\item We evaluate DANTE against several baselines within a high fidelity SUMO driven simulation testbed that reproduces real world urban road networks across six locations, spanning dense urban grids, high speed highway corridors, coastal boulevards, and campus road networks, under realistic time varying mobility and V2X link dynamics.
\end{enumerate}

\section{Related Work}
\label{sec:related_work}

\begin{table*}[t]
\centering
\caption{Comparison of DANTE with representative baselines across the key design dimensions of decentralized personalized learning in vehicular networks.}
\label{tab:comparison}
\renewcommand{\arraystretch}{1.15}
\small
\begin{tabular}{l c c c c c c}
\hline
\textbf{Method}
  & \textbf{Adaptive}
  & \textbf{Personalized}
  & \textbf{Resource}
  & \textbf{Dual}
  & \textbf{Byzantine}
  & \textbf{Learned} \\
  & \textbf{Topology}
  &
  & \textbf{Aware}
  & \textbf{Interface}
  & \textbf{Robust}
  & \textbf{Selection} \\
\hline
FedAvg~\cite{mcmahan2017communication}
  & \texttimes & \texttimes & \texttimes
  & \texttimes & \texttimes & \texttimes \\
D-PSGD~\cite{lian2017dsgd}
  & \texttimes & \texttimes & \texttimes
  & \texttimes & \texttimes & \texttimes \\
CHOCO-SGD~\cite{koloskova2019choco}
  & \texttimes & \texttimes & \checkmark
  & \texttimes & \texttimes & \texttimes \\
BALANCE~\cite{fang2024balance}
  & \texttimes & \texttimes & \texttimes
  & \texttimes & \checkmark & \texttimes \\
MDFL~\cite{chen2025mobility}
  & \checkmark & \texttimes & \checkmark
  & \texttimes & \texttimes & \checkmark \\
CDPL~\cite{boubouh2020robust}
  & \checkmark & \checkmark & \texttimes
  & \texttimes & \checkmark & \checkmark \\
DPFL~\cite{kharrat2025dpfl}
  & \checkmark & \checkmark & \checkmark
  & \texttimes & \texttimes & \texttimes \\
pFedGraph~\cite{ye2023pfedgraph}
  & \checkmark & \checkmark & \texttimes
  & \texttimes & \texttimes & \texttimes \\

\hline
\textbf{DANTE (ours)}
  & \checkmark & \checkmark & \checkmark
  & \checkmark & \checkmark & \checkmark \\
\hline
\end{tabular}
\end{table*}


The growing demand for privacy-preserving, low-latency collaborative intelligence at the network edge has made decentralized learning a major research direction, with a large body of work now addressing serverless optimization, personalized collaboration, communication efficiency, and topology design. Early decentralized federated optimization focused on replacing server-based aggregation with peer-to-peer mixing. D-PSGD~\cite{lian2017dsgd} established the canonical gossip-based training paradigm in which clients alternate local SGD and neighbor averaging over a fixed mixing graph. D$^2$~\cite{tang2018d2} strengthened this line by removing the dominant penalty induced by data heterogeneity, while SGP~\cite{assran2019sgp} extended decentralized training to directed and time-varying communication graphs through PushSum-based gossip. CHOCO-SGD~\cite{koloskova2019choco} further reduced communication overhead by combining decentralized optimization with compressed message exchange. These methods provide the algorithmic foundation of decentralized federated learning, but they still optimize a shared objective over generic peer graphs and therefore do not address personalized collaboration.

This limitation becomes particularly restrictive in vehicular networks, where clients are mobile, communication opportunities are short-lived, and local data distributions are strongly context dependent. Barbieri et al.~\cite{barbieri2022extended} demonstrated that consensus-driven decentralized federated learning can support extended sensing in connected vehicles through direct V2V model exchange. Yuan et al.~\cite{yuan2023p2p} proposed peer-to-peer federated continual learning for naturalistic driving action recognition, showing that decentralized exchange can preserve privacy while mitigating catastrophic forgetting. Chen et al.~\cite{chen2025mobility} moved closer to practical V2X deployment by optimizing local iteration counts and leader selection under mobility and wireless constraints through multi-agent reinforcement learning. Although these studies bring decentralized learning into vehicular settings, they primarily target efficient coordination of shared-model training and do not jointly address personalization, dual link mode for efficient and low latency communication, and Byzantine-aware collaboration.

Because vehicular data are inherently non-IID, personalized decentralized learning is a more natural formulation than consensus on a single global model. Various works~\cite{vanhaesebrouck2017decentralized,boubouh2020robust,basmadjian2022advantages} explored the concept of personalized learning over peer-to-peer networks. \cite{boubouh2020robust} proposed CDPL, a robust and adaptive algorithm to learn personalized models in dynamic networks by evaluating the received updates against a small subset of the peer's local dataset. The authors in~\cite{zantedeschi2020fully} suggested to jointly learn personalized models and sparse collaboration graphs in a fully decentralized setting. DisPFL~\cite{dai2022dispfl} introduced decentralized sparse training and personalized masks, reducing both communication and computation while maintaining client-specific models. DPFL~\cite{kharrat2025dpfl}, which we also implement, cast decentralized personalization as a bilevel optimization problem and used constrained greedy graph construction to select collaborators under a fixed budget. These methods establish that who collaborates with whom is central to personalization quality, yet they do not account for the rapidly changing neighborhoods, asymmetric interface costs, and trust uncertainty that define vehicular V2X environments.

Operational constraints add a second layer of difficulty. Communication and computation efficiency have been studied through compression, resource allocation, and graph sparsification. CHOCO-SGD~\cite{koloskova2019choco} reduces payload size through compressed gossip, while Meng et al.~\cite{meng2025resource} jointly optimize neighbor selection, transmit power, and CPU frequency using graph neural networks in decentralized wireless learning. Zhang et al.~\cite{zhang2024energy} showed that graph sparsification can substantially reduce energy consumption at bottleneck nodes while preserving learning quality. Closely related to these efforts is the topology-optimization literature, which treats the collaboration graph itself as a statistical design variable. pFedGraph~\cite{ye2023pfedgraph} infers collaboration weights from model similarity and dataset size to produce client-specific aggregation patterns. DFedPGP~\cite{liu2024dfedpgp} showed that directed collaboration can outperform symmetric peer exchange in personalized decentralized learning. Le Bars et al.~\cite{lebars2023topology} further demonstrated that topology should be adapted to data heterogeneity itself, and formalized topology learning as an optimization problem driven by neighborhood heterogeneity. Together, these works show that resource allocation and topology design are decisive for decentralized learning performance, but they stop short of unifying personalization, hybrid V2X link heterogeneity, and adversarial robustness in a single decentralized control framework.

DANTE differs from all of these approaches in three coupled ways. First, it formulates vehicular collaboration as a dynamic network formation problem rather than fixed gossip, centralized graph inference, or static budgeted graph construction. Second, it explicitly differentiates the cost and reachability of PC5 sidelink and Uu relay communication, allowing each vehicle to trade off local low-cost exchange against longer-range but more expensive collaboration. Third, it embeds trust estimation and robust aggregation directly into topology control, so collaborator selection depends jointly on personalized learning value, communication cost, and reliability. To the best of our knowledge, no prior method combines these properties in a fully decentralized personalized learning framework for hybrid 5G C-V2X networks.



\section{System Model and Problem Definition}
\label{sec:system_model}

% We consider a fully decentralized vehicular learning system operating over discrete decision rounds $t=0,1,\dots,T-1$. Let $\mathcal{V}_t$ denote the set of active vehicles at round $t$. Due to mobility, the set of vehicles and their communication opportunities vary over time.

We consider a fully decentralized vehicular learning system in which connected vehicles collaboratively train personalized models over a hybrid 5G C-V2X network. Each vehicle aims to learn a model adapted to its own local context while selectively incorporating updates from statistically compatible vehicles. Collaboration is performed without any central coordination and relies on both PC5 sidelink and 5G Uu communication links, which exhibit different range, latency, and cost characteristics. In the following, we formalize the network topology enabling personalized training and its associated computation and communication cost model.

\subsection{Collaboration Graph}

We consider a time-varying vehicular network observed over discrete decision rounds $t \in \{0,\dots,T-1\}$. At each round $t$, let $\mathcal{V}_t$ denote the set of active vehicles participating in collaborative learning. Since vehicle locations and wireless connectivity evolve continuously, the communication topology is inherently dynamic. Model exchange is supported by a hybrid V2X network with two communication modes: direct PC5 sidelink and infrastructure-assisted 5G Uu relay. We represent feasible communication opportunities by the graph
%
\begin{equation*}
    \mathcal{G}_t^{\mathrm{feas}} = \bigl(\mathcal{V}_t,\mathcal{E}_t^{\mathrm{sl}},\mathcal{E}_t^{\mathrm{uu}}\bigr),
\end{equation*}
%
where $(i,j)\in\mathcal{E}_t^{\mathrm{sl}}$ if a direct PC5 transmission between vehicles $i$ and $j$ is feasible at round $t$, and $(i,j)\in\mathcal{E}_t^{\mathrm{uu}}$ if a 5G Uu path through the infrastructure is available. For each pair $(i,j)$, we define the feasible communication mode set as
%
\begin{equation*}
    \mathcal{M}_{ij}^{(t)} = \{m\in\{\mathrm{sl},\mathrm{uu}\}:(i,j)\in\mathcal{E}_t^{m}\}.
\end{equation*}
%
The feasible neighbor set of vehicle $i$ is then
\begin{equation*}
    \mathcal{F}_i^{(t)} = \{j\in\mathcal{V}_t\setminus\{i\}:\mathcal{M}_{ij}^{(t)}\neq\emptyset\}.
\end{equation*}
%
The actual collaboration is modeled by the directed graph
%
\begin{equation*}
    \mathcal{G}_t^{\mathrm{coll}} = \bigl(\mathcal{V}_t,\mathcal{E}_t^{\mathrm{coll}}\bigr),
\end{equation*}
%
where $(j\!\to\! i)\in\mathcal{E}_t^{\mathrm{coll}}$ indicates that vehicle $i$ uses the update received from vehicle $j$ at round $t$. The collaboration neighborhood of vehicle $i$ is therefore
%
\begin{equation*}
    C_i^{(t)} = \{j\in\mathcal{F}_i^{(t)}:(j\!\to\! i)\in\mathcal{E}_t^{\mathrm{coll}}\}.
\end{equation*}
%
For each active collaboration edge, vehicle $i$ selects a communication mode $m_{ij}^{(t)}\in\mathcal{M}_{ij}^{(t)}$ and assigns a collaboration weight $\alpha_{ij}^{(t)}\geq 0$, subject to
%
\begin{equation*}
    \sum_{j\in C_i^{(t)}} \alpha_{ij}^{(t)} \leq 1,
    \qquad
    \alpha_{ii}^{(t)} = 1 - \sum_{j\in C_i^{(t)}} \alpha_{ij}^{(t)}.
\end{equation*}
%
Hence, vehicle $i$ always retains a fraction of its own model while selectively incorporating updates from its collaborators. The collaboration graph is asymmetric and overlapping, that is, $j\in C_i^{(t)}$ does not imply $i\in C_j^{(t)}$.


%% ----------------------------------------------------------------

\subsection{Learning Model}
\label{subsec:learning}

Each vehicle $i\in\mathcal{V}_t$ holds a private dataset $\mathcal{D}_i=\{(x_{i,k},y_{i,k})\}_{k=1}^{n_i}$ drawn from a local distribution $P_i$. The distributions $\{P_i\}$ are heterogeneous across vehicles, yielding non-IID data. Vehicle $i$ maintains a personalized model $w_i^{(t)}\in\mathbb{R}^d$ and aims to minimize its population risk
%
\begin{equation*}
    F_i(w)=\mathbb{E}_{(x,y)\sim P_i}\bigl[\ell(w;x,y)\bigr].
\end{equation*}
%
In practice, since $P_i$ is unknown, local optimization is performed on the empirical risk
%
\begin{equation*}
    \hat{F}_i(w)=\frac{1}{n_i}\sum_{k=1}^{n_i}\ell(w;x_{i,k},y_{i,k}).
\end{equation*}
%
The personalized objective of vehicle $i$ is therefore
\begin{equation*}
    \min_{w_i} F_i(w_i).
\end{equation*}
%
Since $P_i\neq P_j$ in general, enforcing $w_i=w_j$ for all vehicles is suboptimal. At the system level, the objective is to learn the set of personalized models $W=\{w_i\}_{i\in\mathcal{V}_t}$ by minimizing the aggregate multi-task loss
\begin{equation*}
    \min_W \mathcal{L}(W) = \sum_{i\in\mathcal{V}_t}\rho_i F_i(w_i),
    \qquad
    \rho_i\geq 0,
    \quad
    \sum_{i\in\mathcal{V}_t}\rho_i=1,
\end{equation*}
%
where $\rho_i$ denotes the weight of vehicle $i$. Unlike classical Decentralized Learning (DL), this objective does not impose model consensus.

% At the beginning of round $t$, vehicle $i$ aggregates its model with the models received from its collaborators:
At the beginning of round $t$, vehicle $i$ receives the most recent models $\hat{w}_j^{(t)}$ from its collaborators $j\in C_i^{(t)}$ and forms the aggregate model $\bar{w}_i^{(t)}$.
%
\begin{equation}
    \bar{w}_i^{(t)} = \alpha_{ii}^{(t)} w_i^{(t)} + \sum_{j\in C_i^{(t)}} \alpha_{ij}^{(t)} \hat{w}_j^{(t)},
\end{equation}
%
where $\hat{w}_j^{(t)}$ is the model received from vehicle $j$, and
%
\begin{equation*}
    \alpha_{ii}^{(t)}+\sum_{j\in C_i^{(t)}}\alpha_{ij}^{(t)}=1,
    \qquad
    \alpha_{ij}^{(t)}\geq 0.
\end{equation*}
%
Starting from $\bar{w}_i^{(t)}$, vehicle $i$ performs $I_i^{(t)}$ local SGD steps on mini-batches of size $b_i$ sampled from $\mathcal{D}_i$:
%
\begin{equation}
    w_i^{(t,\ell+1)} = w_i^{(t,\ell)} - \eta_i^{(t)} g_i\!\left(w_i^{(t,\ell)};\xi_i^{(t,\ell)}\right),
\end{equation}
%
for $\ell=0,\dots,I_i^{(t)}-1$, with $w_i^{(t,0)}=\bar{w}_i^{(t)}$, where $\xi_i^{(t,\ell)}$ is a mini-batch of size $b_i$ drawn from $\mathcal{D}_i$ and $g_i(\cdot;\xi_i^{(t,\ell)})$ is the corresponding stochastic gradient of $\hat{F}_i$. The updated personalized model is then
%
\begin{equation*}
    w_i^{(t+1)} = w_i^{(t,I_i^{(t)})}.
\end{equation*}

%% ----------------------------------------------------------------

% \subsection{Assumptions and constraints.}
% We assume that each local objective $F_i$ is differentiable and $L$-smooth, and that stochastic gradients are unbiased with bounded variance. Each vehicle has access only to local observations and does not observe the global graph or the states of non-neighboring vehicles. Actions must satisfy the feasibility and budget constraints
% \begin{align}
% &C_i(t)\subseteq \mathcal{F}_i(t), \qquad |C_i(t)|\leq K_i,\\
% &\sum_{j\in C_i(t)} E_{ij}^{m_{ij}(t)}(t)+E_i^{\mathrm{cmp}}(t)\leq \bar{E}_i,\\
% &\sum_{j\in C_i(t)} B_{ij}^{m_{ij}(t)}(t)\leq \bar{B}_i,\\
% &\tau_i^{\mathrm{cmp}}(t)+\sum_{j\in C_i(t)} \tau_{ij}^{m_{ij}(t)}(t)\leq \bar{\tau}_i,
% \end{align}
% where $\bar{E}_i$, $\bar{B}_i$, and $\bar{\tau}_i$ denote the per-round energy, bandwidth, and latency budgets of vehicle $i$.

%% ----------------------------------------------------------------

\subsection{Computation and Communication Model}

Each vehicle incurs both local computation cost and V2X communication cost during collaborative training. We model these costs at the granularity of
one learning round.

If vehicle $i$ performs $I_i^{(t)}$ local SGD steps on mini-batches of size $b_i$, with $c_i$ CPU cycles per sample and processor frequency
$f_i^{(t)}$, its local computation time and energy are
%
\begin{equation*}
    \tau_{i,\mathrm{cmp}}^{(t)} = \frac{I_i^{(t)}\, b_i\, c_i}{f_i^{(t)}},
    \qquad
    E_{i,\mathrm{cmp}}^{(t)} = \kappa I_i^{(t)}\, b_i\, c_i \bigl(f_i^{(t)}\bigr)^2
\end{equation*}
%
where $\kappa$ is the effective switched capacitance.

Vehicles exchange model updates over hybrid 5G C-V2X links, either through direct PC5 sidelink or infrastructure-assisted 5G Uu relay, which exhibit
different range, latency, and power characteristics \cite{garcia2021tutorial,molina2017ltev,3gpp2021architecture}. Let $S$ denote the model payload size in bits. For a PC5 sidelink transmission from vehicle $j$ to vehicle $i$, the achievable rate is
%
\begin{equation*}
    R_{ji,\mathrm{sl}}^{(t)} = B_{ji,\mathrm{sl}}^{(t)} \log_2\!\bigl(1+\mathrm{SINR}_{ji,\mathrm{sl}}^{(t)}\bigr),
\end{equation*}
%
with transmission time and energy
%
\begin{equation*}
    \tau_{ji,\mathrm{sl}}^{(t)}=\frac{S}{R_{ji,\mathrm{sl}}^{(t)}},
    \qquad
    E_{ji,\mathrm{sl}}^{(t)}=p_{j,\mathrm{sl}}^{(t)}\tau_{ji,\mathrm{sl}}^{(t)}.
\end{equation*}
%
For a 5G Uu relay path through the serving gNB, the end-to-end transfer time is determined by the uplink and downlink legs:
\begin{equation*}
    \tau_{ji,\mathrm{uu}}^{(t)} = \frac{S}{R_{j\rightarrow g}^{u,(t)}} + \frac{S}{R_{g\rightarrow i}^{d,(t)}},
\end{equation*}
%
and the associated communication energy is
%
\begin{equation*}
    E_{ji,\mathrm{uu}}^{(t)} = p_j^{u,(t)}\frac{S}{R_{j\rightarrow g}^{u,(t)}} + p_i^{\mathrm{rx},(t)}\frac{S}{R_{g\rightarrow i}^{d,(t)}}.
\end{equation*}
%
This model captures the core asymmetry of hybrid V2X collaboration: PC5 is typically cheaper and lower-latency but range-limited, whereas 5G Uu relay
offers larger reach at higher communication cost.


%% ----------------------------------------------------------------

\subsection{Threat Model}

This work considers an open vehicular environment in which a subset $\mathcal{B}_t\subseteq\mathcal{V}_t$ may behave in a Byzantine manner at round $t$. A Byzantine vehicle is assumed to know the protocol, collude with other Byzantine vehicles, and send arbitrary, inconsistent, or strategically poisoned model updates, but it cannot directly compromise the local computation of honest vehicles. The threat model covers crash and communication failures, free-riding or lazy participants that consume updates without contributing useful models, and adversarial model poisoning aimed at introducing backdoors, causing divergence, or steering learning toward ineffective solutions. We particularly consider the case where Byzantine vehicles collude to dominate the neighborhood of an honest vehicle and corrupt its personalized model. To mitigate the effect of unreliable or Byzantine neighbors, each vehicle maintains a local trust score $q_{ij}^{(t)}\in[0,1]$ for every feasible neighbor $j\in\mathcal{F}_i^{(t)}$. This score represents vehicle $i$'s current estimate of the reliability of $j$'s transmitted model updates.



\section{DANTE: Dynamic Network Formation and Decentralized Solution}
\label{sec:dante}

Having defined the system model in Section~\ref{sec:system_model}, we now answer the central question of this work: \emph{how should each vehicle decide, at every round, with whom to collaborate so as to maximize its personalized learning benefit under the energy, bandwidth, latency, and trust constraints of a hybrid 5G C-V2X network?} We proceed in two steps. First, we cast vehicular decentralized personalized learning (VDPL) as a \emph{dynamic network formation game} played over the time-varying collaboration graph, in which each vehicle is a strategic player that forms and severs directed links to its peers. Second, we introduce \textsc{Dante}~(\underline{D}ual-interface \underline{A}ttention-guided \underline{N}etwork-formation for \underline{T}rustworthy vehicular \underline{E}dge learning), a fully decentralized solver in which every vehicle learns its own link-formation policy via an attention-guided PPO actor--critic, aggregates the received models through a fused attention--trust--robustness weighting, and maintains a local trust-update loop. No central critic, parameter server, or globally observed state is required at any stage.

%-----------------------------------------------------------
\subsection{Dynamic Network Formation as a Strategic Markov Game}
\label{subsec:game}

We model collaboration in VDPL as a discrete-time strategic network formation game in the tradition of Jackson and Wolinsky~\cite{jackson1996strategic,bala2000noncooperative,jackson2002evolution}. This class of models treats agents as deciding which directed links to form or sever and lets payoffs depend on the resulting graph rather than on a partition of the agent set into disjoint coalitions. This graph-level abstraction is well suited to vehicular collaboration, where links are overlapping, directed, and asymmetric: vehicle $i$ may benefit from $j$'s update without any reciprocity obligation, and each vehicle observes its environment only locally.

\begin{definition}[Dynamic network formation game for VDPL]
\label{def:game}
A dynamic network formation game for VDPL is the tuple
\begin{equation*}
    \mathfrak{G} = \Bigl(\,\{\mathcal{V}_t\}_{t\ge 0},\;\mathcal{S},\;\{\mathcal{A}_i(\cdot)\}_i,\;\mathcal{P},\;\{u_i\}_i,\;\gamma\,\Bigr),
\end{equation*}
whose components are defined as follows.
\begin{itemize}
    \item \textbf{Players.} At round $t$, the players are the active vehicles $\mathcal{V}_t$ of Section~\ref{sec:system_model}. Every honest vehicle is a strategic, self-interested agent that maximizes its own personalized learning performance under local resource constraints.

    \item \textbf{State and local observation.} The global state $s^{(t)}\in\mathcal{S}$ collects vehicle positions, channel conditions, resource budgets $b_i^{(t)}$, personalized models $\{w_i^{(t)}\}$, and trust vectors $\{q_{ij}^{(t)}\}$. No vehicle observes $s^{(t)}$. Vehicle $i$ receives only a local observation $o_i^{(t)}\in\mathcal{O}_i$ that exposes its own state together with per-candidate link features $\{f_{ij}^{(t)}\}$ and trust scores $\{q_{ij}^{(t)}\}$ restricted to the feasible neighbor set $\mathcal{F}_i^{(t)}$ defined in Section~\ref{sec:system_model}.

    \item \textbf{Feasible actions.} Vehicle $i$ plays a per-candidate binary vector
    \begin{equation*}
        a_i^{(t)} = \bigl(a_{ij}^{(t)}\bigr)_{j\in\mathcal{F}_i^{(t)}}\in\{0,1\}^{|\mathcal{F}_i^{(t)}|},
    \end{equation*}
    where $a_{ij}^{(t)}=1$ accepts the model update of neighbor $j$ at round $t$. Actions induce the collaboration set $C_i^{(t)}=\{j\in\mathcal{F}_i^{(t)}:a_{ij}^{(t)}=1\}$ and the directed collaboration graph $\mathcal{G}_t^{\mathrm{coll}}$ of Section~\ref{sec:system_model}, with $(j\!\to\! i)\in\mathcal{E}_t^{\mathrm{coll}}\iff a_{ij}^{(t)}=1$. Since every candidate $j\in\mathcal{F}_i^{(t)}$ is already tagged with its physical interface $m_{ij}^{(t)}\in\mathcal{M}_{ij}^{(t)}$ by the communication layer, vehicle $i$ chooses directed edges rather than separately optimizing the interface of a fixed peer.

    \item \textbf{Per-round payoff.} Vehicle $i$ receives the instantaneous utility
    \begin{equation}
    \label{eq:stage_payoff}
        u_i^{(t)} = \Delta_i^{(t)} - \lambda_E\,E_i^{(t)} - \lambda_B\,B_i^{(t)} - \lambda_T\,T_i^{(t)},
    \end{equation}
    where
    \begin{equation*}
        \Delta_i^{(t)} = \widehat{F}_{i,\mathrm{val}}\!\bigl(w_i^{(t)}\bigr) - \widehat{F}_{i,\mathrm{val}}\!\bigl(w_i^{(t+1)}\bigr)
    \end{equation*}
    is the loss reduction measured on a held-out local validation split $\mathcal{D}_i^{\mathrm{val}}\subset\mathcal{D}_i$, $E_i^{(t)}$, $B_i^{(t)}$, and $T_i^{(t)}$ are the round-level energy, bandwidth, and latency costs inherited from the computation and communication model of Section~\ref{sec:system_model}, and $\lambda_E,\lambda_B,\lambda_T\ge 0$ are per-vehicle cost weights.

    \item \textbf{Transitions.} The state evolves according to a kernel $\mathcal{P}(\cdot\mid s^{(t)},a^{(t)})$ jointly induced by vehicle mobility, wireless-channel stochasticity, resource depletion $b_i^{(t+1)}=b_i^{(t)}-[E_i^{(t)},B_i^{(t)}]^\top$, trust evolution (Section~\ref{subsubsec:trust}), and the personalized model update (Section~\ref{subsec:dante}). Vehicle $i$ does not observe $\mathcal{P}$ and interacts with it only through $o_i^{(t)}$.
\end{itemize}
\end{definition}

\noindent\textbf{Long-term objective.} Given a joint policy profile $\pi=(\pi_i)_i$, vehicle $i$ seeks a local policy that maximizes its discounted infinite-horizon return
\begin{equation}
\label{eq:long_term}
    J_i(\pi_i;\pi_{-i}) = \mathbb{E}_{\pi}\!\left[\sum_{t=0}^{\infty}\gamma^{t}\,u_i^{(t)}\right],\qquad \gamma\in(0,1),
\end{equation}
where $\pi_{-i}$ denotes the policy profile of the remaining vehicles. Because every vehicle optimizes its own return conditional on the strategies of the others, Equation~\eqref{eq:long_term} defines a Markov game over the evolving collaboration graph rather than a joint cooperative objective; we return to the Nash-equilibrium question in Section~\ref{subsec:convergence}.

\noindent\textbf{Constraints.} At each round, the decision of vehicle $i$ must additionally satisfy
\begin{align}
\label{eq:cons_k}
    |C_i^{(t)}| &\le K_i,\\
\label{eq:cons_energy}
    E_{i,\mathrm{cmp}}^{(t)}+\sum_{j\in C_i^{(t)}} E_{ij}^{m_{ij},(t)} &\le \bar E_i^{(t)},\\
\label{eq:cons_bw}
    \sum_{j\in C_i^{(t)}} B_{ij}^{m_{ij},(t)} &\le \bar B_i^{(t)},\\
\label{eq:cons_lat}
    T_{i,\mathrm{cmp}}^{(t)}+\max_{j\in C_i^{(t)}} T_{ij}^{m_{ij},(t)} &\le \bar T_i^{(t)},
\end{align}
where $K_i$ is the per-vehicle neighborhood cap and $\bar E_i^{(t)},\bar B_i^{(t)},\bar T_i^{(t)}$ are the per-round energy, bandwidth, and latency budgets of Section~\ref{sec:system_model}. Latency~\eqref{eq:cons_lat} and neighborhood size~\eqref{eq:cons_k} are hard per-round caps enforced at action-admissibility time, whereas the energy and bandwidth caps~\eqref{eq:cons_energy}--\eqref{eq:cons_bw} are tracked through the budget state $b_i^{(t)}$ so that a vehicle that overspends early loses collaboration capacity later. This design keeps the problem amenable to policy-gradient methods without resorting to a Lagrangian relaxation, while still penalizing budget depletion endogenously through the return.

%-----------------------------------------------------------
\subsection{Utility of Collaboration}
\label{subsec:utility}

The stage payoff~\eqref{eq:stage_payoff} uses the validation-loss drop $\Delta_i^{(t)}$ as its observable learning signal. That quantity is suitable for online decision making but only reveals the value of collaboration after a round has been completed. To expose the learning geometry underlying $\Delta_i^{(t)}$ and to support the theoretical analysis, we introduce an analytical surrogate based on gradient agreement.

Let $g_i^{(t)}$ and $g_j^{(t)}$ denote stochastic gradients of vehicles $i$ and $j$ evaluated at their current iterates. We measure directional agreement by the cosine similarity
\begin{equation*}
    s_{ij}^{(t)} \;=\; \frac{\bigl(g_i^{(t)}\bigr)^\top g_j^{(t)}}{\|g_i^{(t)}\|\,\|g_j^{(t)}\|+\varepsilon}\;\in\;[-1,1],
\end{equation*}
with $\varepsilon>0$ a numerical-stability constant. This quantity admits a first-order interpretation: Taylor-expanding $F_i$ around $w_i^{(t)}$ gives
\begin{equation*}
    F_i\!\bigl(w_i^{(t)}-\eta\, g_j^{(t)}\bigr) = F_i\!\bigl(w_i^{(t)}\bigr) - \eta\,\bigl(\nabla F_i(w_i^{(t)})\bigr)^{\!\top} g_j^{(t)} + O(\eta^2),
\end{equation*}
so under the standard assumption that $g_i^{(t)}$ is unbiased for $\nabla F_i(w_i^{(t)})$, a positive $s_{ij}^{(t)}$ indicates, in expectation, that $j$'s update direction is aligned with the local descent direction of $i$. Combining alignment with the trust score $q_{ij}^{(t)}\in[0,1]$ introduced in Section~\ref{subsubsec:trust}, we define the trust-weighted collaboration surrogate
\begin{equation}
\label{eq:G_surrogate}
    G_i\!\bigl(C_i^{(t)},s^{(t)}\bigr) \;=\; \sum_{j\in C_i^{(t)}} q_{ij}^{(t)}\,s_{ij}^{(t)}.
\end{equation}
Large values of $G_i$ correspond to collaboration sets whose update directions are, on average, both aligned with the local descent of $i$ and originating from trusted peers, whereas negative values indicate that the selected neighbors tend to push the model away from descent. The surrogate is introduced for interpretation and convergence analysis (Section~\ref{subsec:convergence}); the reward that drives the policy is the observable $\Delta_i^{(t)}$ entering~\eqref{eq:stage_payoff}.

%-----------------------------------------------------------
\subsection{\textsc{Dante}: Attention-Guided Decentralized Solver}
\label{subsec:dante}

\textsc{Dante} factorizes the per-vehicle policy through three coupled components: (i) a graph-attention policy head that scores candidate peers from local link features, (ii) a robust personalized aggregation layer that fuses attention with trust and a coordinate-wise robustness filter, and (iii) a local trust update that closes the loop by down-weighting peers whose past contributions hurt the local loss. Every component runs locally on vehicle $i$ using $o_i^{(t)}$ alone; no joint state, joint critic, or shared parameters are ever assembled.

\subsubsection{GAT policy parameterization}
\label{subsubsec:gatppo}

Let $x_i^{(t)}\in\mathbb{R}^{d_o}$ be the ego feature vector of vehicle $i$ and $x_{ij}^{(t)}\in\mathbb{R}^{d_n}$ the feature vector of candidate neighbor $j\in\mathcal{F}_i^{(t)}$. A single-head graph-attention encoder computes
\begin{equation}
\label{eq:gat}
\begin{split}
    h_i^{(t)} &= W_s\, x_i^{(t)},\qquad v_{ij}^{(t)} = W_n\, x_{ij}^{(t)},\\
    \beta_{ij}^{(t)} &= \frac{\exp\bigl(e_{ij}^{(t)}\bigr)}{\sum_{k\in\mathcal{F}_i^{(t)}}\exp\bigl(e_{ik}^{(t)}\bigr)},
\end{split}
\end{equation}
with unnormalized logits $e_{ij}^{(t)}=a^{\!\top}\mathrm{LeakyReLU}\!\bigl([h_i^{(t)}\,\|\,v_{ij}^{(t)}]\bigr)$ and learnable parameters $W_s$, $W_n$, and $a$. The neighborhood context is $c_i^{(t)}=\sum_{j\in\mathcal{F}_i^{(t)}}\beta_{ij}^{(t)}v_{ij}^{(t)}$. The raw attention $\beta_{ij}^{(t)}$ plays a dual role: it drives the actor's Bernoulli logits, and, combined with the trust and robustness scores introduced below, it shapes the final aggregation weights $\alpha_{ij}^{(t)}$. This coupling ties \emph{whom the policy chooses to listen to} directly to \emph{how much weight that peer carries in the model update}.

Conditioned on $(h_i^{(t)},c_i^{(t)},v_{ij}^{(t)})$, the actor outputs a per-candidate Bernoulli and the critic outputs a scalar value:
\begin{equation}
\label{eq:actor_critic}
\begin{split}
    \pi_{\theta_i}\!\bigl(a_{ij}^{(t)}\!=\!1\mid o_i^{(t)}\bigr) &= \sigma\!\Bigl(f_{\theta_i}\!\bigl([h_i^{(t)}\,\|\,c_i^{(t)}\,\|\,v_{ij}^{(t)}]\bigr)\Bigr),\\
    V_{\phi_i}\!\bigl(o_i^{(t)}\bigr) &= g_{\phi_i}\!\bigl([h_i^{(t)}\,\|\,c_i^{(t)}]\bigr),
\end{split}
\end{equation}
where $\sigma$ is the logistic sigmoid. Because the per-candidate actions are conditionally independent Bernoullis, the joint log-probability factorizes as $\log\pi_{\theta_i}(a_i^{(t)}\mid o_i^{(t)})=\sum_{j\in\mathcal{F}_i^{(t)}}\log\pi_{\theta_i}(a_{ij}^{(t)}\mid o_i^{(t)})$, making per-vehicle PPO updates tractable regardless of $|\mathcal{F}_i^{(t)}|$. Each vehicle stores only its own parameters $(\theta_i,\phi_i)$: no critic ever observes a joint state or joint action, and the scheme avoids the realizability issues of centralized-training-with-decentralized-execution schemes in V2X, where a joint state would presuppose a central entity that our setting precludes. If more than $K_i$ candidates are sampled as active, \textsc{Dante} retains the $K_i$ highest-attention edges, which enforces the neighborhood cap~\eqref{eq:cons_k} by construction.

\subsubsection{Trust score and robustness filter}
\label{subsubsec:trust}

Vehicle $i$ maintains a local trust score $q_{ij}^{(t)}\in[0,1]$ for every feasible neighbor $j\in\mathcal{F}_i^{(t)}$, interpreted as vehicle $i$'s current estimate of the reliability of $j$'s transmitted updates. The score is updated by an exponentially weighted recursion
\begin{equation}
\label{eq:trust_update}
    q_{ij}^{(t+1)} = (1-\rho)\,q_{ij}^{(t)} + \rho\,\phi_{ij}^{(t)},\qquad \rho\in(0,1],
\end{equation}
where the per-round reliability signal $\phi_{ij}^{(t)}\in[0,1]$ is defined as the product of two local indicators,
\begin{equation*}
    \phi_{ij}^{(t)} = \mathbf{1}\!\bigl\{r_{ij}^{(t)}\ge r_{\min}\bigr\}\cdot\mathbf{1}\!\bigl\{\widehat F_{i,\mathrm{val}}(w_i^{(t+1)})\le \widehat F_{i,\mathrm{val}}(w_i^{(t)})+\tau_{\mathrm{val}}\bigr\},
\end{equation*}
namely (i) vehicle $j$ passed the local robust filter at round $t$ and (ii) the validation loss of vehicle $i$ did not increase after including $j$'s update, up to a small slack $\tau_{\mathrm{val}}\ge 0$. Because $\phi_{ij}^{(t)}\in[0,1]$ and $\rho\in(0,1]$, the recursion keeps $q_{ij}^{(t)}\in[0,1]$ for every $t$. The smoothing in~\eqref{eq:trust_update} gives persistence against transient link errors while still allowing a sustained Byzantine peer to be progressively de-weighted.

The robustness score $r_{ij}^{(t)}\in[0,1]$ is produced by a coordinate-wise robust filter applied to the set of received models $\{\hat w_j^{(t)}\}_{j\in C_i^{(t)}}$. Let $\mathrm{TM}_i^{(t)}$ denote the coordinate-wise trimmed mean of $\{\hat w_j^{(t)}\}_{j\in C_i^{(t)}}$ and let $d_{ij}^{(t)}=\|\hat w_j^{(t)}-\mathrm{TM}_i^{(t)}\|_2$; we set $r_{ij}^{(t)}=\exp\!\bigl(-d_{ij}^{(t)}/\sigma_r\bigr)\in(0,1]$, so that peers that drift far from the trimmed-mean consensus contribute near-zero weight while still-close peers retain near-unit score. Coordinate-wise trimmed mean is a standard robust aggregator in federated learning and tolerates up to a fraction $f/|C_i^{(t)}|<1/2$ of Byzantine neighbors with bounded honest gradients~\cite{yin2018byzantine,karimireddy2021learning}.

\subsubsection{Fused attention--trust--robustness aggregation}
\label{subsubsec:agg}

\textsc{Dante} combines the GAT attention $\beta_{ij}^{(t)}$, the trust score $q_{ij}^{(t)}$, and the robustness score $r_{ij}^{(t)}$ into the final aggregation weight that instantiates the collaborative update of Section~\ref{subsec:learning}:
\begin{equation}
\label{eq:alpha_fused}
    \alpha_{ij}^{(t)} = \bigl(1-\alpha_{ii}^{(t)}\bigr)\cdot\frac{\beta_{ij}^{(t)}\,q_{ij}^{(t)}\,r_{ij}^{(t)}}{\varepsilon+\displaystyle\sum_{k\in C_i^{(t)}}\beta_{ik}^{(t)}\,q_{ik}^{(t)}\,r_{ik}^{(t)}},\quad j\in C_i^{(t)},
\end{equation}
with self-retention weight $\alpha_{ii}^{(t)}\in(0,1)$ controlling the personalization/collaboration trade-off. By construction,
\begin{equation*}
    \alpha_{ii}^{(t)}+\sum_{j\in C_i^{(t)}}\alpha_{ij}^{(t)} = 1,\qquad \alpha_{ij}^{(t)}\ge 0,
\end{equation*}
so~\eqref{eq:alpha_fused} satisfies the convex-combination constraint of Section~\ref{subsec:learning} and can be substituted directly into the aggregation rule of that section, yielding $\bar w_i^{(t)}=\alpha_{ii}^{(t)}w_i^{(t)}+\sum_{j\in C_i^{(t)}}\alpha_{ij}^{(t)}\hat w_j^{(t)}$. Starting from $\bar w_i^{(t)}$, vehicle $i$ then executes $I_i^{(t)}$ local SGD steps as prescribed in Section~\ref{subsec:learning} to obtain the next personalized iterate $w_i^{(t+1)}$. The round ordering is therefore \emph{aggregate-then-train}, consistent with the system model.

The fusion in~\eqref{eq:alpha_fused} is also a graceful fallback. If every neighbor in $C_i^{(t)}$ is Byzantine, the robust filter and trust recursion drive $q_{ij}^{(t)}\,r_{ij}^{(t)}\!\to\! 0$, so $\alpha_{ij}^{(t)}\!\to\! 0$ and $\bar w_i^{(t)}\!\to\! w_i^{(t)}$; the collaborative update degenerates to pure local SGD, which is the safest behavior available to an honest vehicle.

\subsubsection{Local PPO optimization}
\label{subsubsec:ppo}

Each vehicle updates its actor--critic pair by Proximal Policy Optimization on local trajectories only. Let $\delta_i^{(t)}=u_i^{(t)}+\gamma\,V_{\phi_i}(o_i^{(t+1)})-V_{\phi_i}(o_i^{(t)})$ denote the temporal-difference error and $\hat A_i^{(t)}=\sum_{h\ge 0}(\gamma\lambda_{\mathrm{GAE}})^h\,\delta_i^{(t+h)}$ the generalized advantage estimate. The clipped PPO surrogate with entropy regularization is
\begin{equation}
\label{eq:ppo}
\begin{split}
    \mathcal{L}_i^{\mathrm{PPO}}(\theta_i) &= \mathbb{E}\!\left[\min\!\bigl(\rho_i^{(t)}\hat A_i^{(t)},\,\mathrm{clip}(\rho_i^{(t)},1-\epsilon,1+\epsilon)\hat A_i^{(t)}\bigr)\right]\\
    &\quad + \eta_{\mathrm{ent}}\,\mathbb{E}\!\bigl[\mathcal{H}\bigl(\pi_{\theta_i}(\cdot\mid o_i^{(t)})\bigr)\bigr],
\end{split}
\end{equation}
with likelihood ratio $\rho_i^{(t)}=\pi_{\theta_i}(a_i^{(t)}\mid o_i^{(t)})/\pi_{\theta_i^{\mathrm{old}}}(a_i^{(t)}\mid o_i^{(t)})$, clip range $\epsilon\in(0,1)$, entropy coefficient $\eta_{\mathrm{ent}}\ge 0$, and Shannon entropy $\mathcal{H}$. The critic regresses on the return targets $\hat R_i^{(t)}=\hat A_i^{(t)}+V_{\phi_i}(o_i^{(t)})$. Each vehicle therefore runs an independent local PPO loop in the style of IPPO~\cite{dewitt2020independent}, without sharing trajectories, critics, or parameters with its peers.

\subsubsection{The DANTE algorithm}
\label{subsubsec:algo}

Algorithm~\ref{alg:dante} summarizes the complete per-vehicle procedure. Each iteration performs one round of environment interaction followed by local learning, and the PPO update is triggered every $K_{\mathrm{upd}}$ rounds from the local rollout buffer. All steps operate on local data and locally exchanged parameters; no global reward, joint state, or joint critic is ever assembled.

\begin{algorithm}[t]
\caption{\textsc{Dante}: per-vehicle decentralized network formation and personalized learning, run locally by each vehicle $i$.}
\label{alg:dante}
\begin{algorithmic}[1]
\Require local dataset $\mathcal{D}_i$; initial model $w_i^{(0)}$; initial trust $q_{ij}^{(0)}\!=\!1$; actor--critic parameters $(\theta_i,\phi_i)$; per-round budgets $(\bar E_i,\bar B_i,\bar T_i)$; neighborhood cap $K_i$; PPO hyperparameters $(\gamma,\lambda_{\mathrm{GAE}},\epsilon,\eta_{\mathrm{ent}},K_{\mathrm{upd}})$; self-weight schedule $\alpha_{ii}^{(t)}$.
\Ensure personalized model $w_i^{(T)}$ and policy $(\theta_i,\phi_i)$.
\For{$t=0,1,\dots,T-1$}
    \State Discover feasible peers $\mathcal{F}_i^{(t)}$ over PC5 sidelink and Uu relay; observe link features $\{f_{ij}^{(t)}\}$ and trust scores $\{q_{ij}^{(t)}\}$.
    \State Encode ego and candidate features $(h_i^{(t)},\{v_{ij}^{(t)}\})$; compute attention $\{\beta_{ij}^{(t)}\}$ and value $V_{\phi_i}(o_i^{(t)})$.\Comment{Eq.~\eqref{eq:gat}}
    \State Sample per-candidate Bernoulli actions $a_{ij}^{(t)}\!\sim\!\pi_{\theta_i}(\cdot\mid o_i^{(t)})$; form $C_i^{(t)}\!=\!\{j:a_{ij}^{(t)}\!=\!1\}$; if $|C_i^{(t)}|\!>\!K_i$, retain the top-$K_i$ edges ranked by $\beta_{ij}^{(t)}$.\Comment{Eq.~\eqref{eq:actor_critic}}
    \State Exchange models with $C_i^{(t)}$ over the selected interface $m_{ij}^{(t)}$; receive $\{\hat w_j^{(t)}\}_{j\in C_i^{(t)}}$.
    \State Apply the coordinate-wise robust filter to obtain $\{r_{ij}^{(t)}\}_{j\in C_i^{(t)}}$.
    \State Form fused weights $\{\alpha_{ij}^{(t)}\}$ and aggregate $\bar w_i^{(t)}=\alpha_{ii}^{(t)}w_i^{(t)}+\sum_{j\in C_i^{(t)}}\alpha_{ij}^{(t)}\hat w_j^{(t)}$.\Comment{Eq.~\eqref{eq:alpha_fused}}
    \State Starting from $\bar w_i^{(t)}$, run $I_i^{(t)}$ local SGD steps to obtain $w_i^{(t+1)}$.\Comment{Section~\ref{subsec:learning}}
    \State Measure $(E_i^{(t)},B_i^{(t)},T_i^{(t)})$ and validation drop $\Delta_i^{(t)}$; form reward $u_i^{(t)}$.\Comment{Eq.~\eqref{eq:stage_payoff}}
    \State Evaluate $\phi_{ij}^{(t)}$ from the robust-filter indicator and the post-update validation loss; update trust $\{q_{ij}^{(t+1)}\}$.\Comment{Eq.~\eqref{eq:trust_update}}
    \State Append transition $(o_i^{(t)},a_i^{(t)},u_i^{(t)},o_i^{(t+1)},V_{\phi_i}(o_i^{(t)}))$ to local rollout buffer $\mathcal{R}_i$.
    \If{$|\mathcal{R}_i|\ge K_{\mathrm{upd}}$}
        \State Compute $\{\hat A_i^{(t)},\hat R_i^{(t)}\}$ via generalized advantage estimation.
        \State Update $\theta_i\!\leftarrow\!\theta_i+\eta_\pi\nabla_{\theta_i}\mathcal{L}_i^{\mathrm{PPO}}(\theta_i)$; update $\phi_i$ by MSE on $\{\hat R_i^{(t)}\}$.\Comment{Eq.~\eqref{eq:ppo}}
        \State Clear $\mathcal{R}_i$.
    \EndIf
\EndFor
\end{algorithmic}
\end{algorithm}

%-----------------------------------------------------------
\subsection{Convergence and Complexity}
\label{subsec:convergence}

The joint problem is non-convex and non-stationary, and global optimality of coupled model and policy updates is out of reach. Following a two-timescale argument standard in multi-agent reinforcement learning, we analyze the learning and policy layers separately: convergence of the personalized learning iterates under a fixed collaboration profile, and stationarity of the local PPO update under an asymptotically stationary local MDP. Detailed proofs are deferred to Appendix~\ref{app:proofs}.

\paragraph{Assumptions.}
We use the following standard conditions.
\begin{enumerate}
    \item[\textbf{A1.}] \emph{(Smoothness)} Each local objective $F_i$ is $L$-smooth and bounded below by $F_i^{\inf}$.
    \item[\textbf{A2.}] \emph{(Unbiased bounded-variance gradients)} $\mathbb{E}[g_i(w;\xi)]=\nabla F_i(w)$ and $\mathbb{E}\|g_i(w;\xi)-\nabla F_i(w)\|^2\le\sigma^2$.
    \item[\textbf{A3.}] \emph{(Robust-aggregation bias)} The fused weights $\{\alpha_{ij}^{(t)}\}$ of~\eqref{eq:alpha_fused} together with the robust filter yield an effective post-aggregation stochastic gradient $\tilde g_i^{(t)}$ satisfying $\|\mathbb{E}[\tilde g_i^{(t)}\mid w_i^{(t)}]-\nabla F_i(w_i^{(t)})\|\le\zeta$ uniformly in $t$, for some $\zeta\ge 0$.
    \item[\textbf{A4.}] \emph{(Bounded rewards and features)} $|u_i^{(t)}|\le u_{\max}$ almost surely, and candidate features and attention outputs are uniformly bounded.
    \item[\textbf{A5.}] \emph{(Persistent exploration)} The entropy bonus in~\eqref{eq:ppo} keeps $\pi_{\theta_i}(\cdot\mid o_i^{(t)})$ strictly positive on every action in the support of an $\epsilon$-perturbation of the current policy.
\end{enumerate}
\textbf{A1}--\textbf{A2} are standard for smooth non-convex SGD~\cite{bottou2018optimization}. \textbf{A3} is the standard interface used in robust federated learning to decouple filter geometry from convergence arguments~\cite{karimireddy2021learning,yin2018byzantine}; the coordinate-wise trimmed mean of Section~\ref{subsubsec:trust} attains $\zeta=\mathcal{O}(f/|C_i^{(t)}|)\cdot B$ for $f$ Byzantine neighbors and bounded honest gradients $\|g_j\|\le B$. \textbf{A4} follows from the bounded validation-loss drop together with the per-round resource caps~\eqref{eq:cons_energy}--\eqref{eq:cons_lat}. \textbf{A5} is enforced by the entropy term in~\eqref{eq:ppo}.

\paragraph{Convergence of the learning layer under fixed policies.}

\begin{proposition}[Robust personalized SGD convergence]
\label{prop:learning}
Fix any joint policy profile $\pi$ for which the induced collaboration sequence $\{C_i^{(t)}\}_{t\ge 0}$ is measurable. Under \textbf{A1}--\textbf{A3}, with diminishing step size $\eta_i^{(t)}=\eta_0/\sqrt{t+1}$ and $\eta_0\le 1/L$, the personalized iterates $\{w_i^{(t)}\}$ generated by aggregate-then-train starting from $\bar w_i^{(t)}$ (Section~\ref{subsec:learning}) satisfy, for every honest vehicle $i$,
\begin{equation}
\label{eq:prop1}
\begin{split}
    \min_{0\le t\le T-1}&\,\mathbb{E}\!\left[\bigl\|\nabla F_i(w_i^{(t)})\bigr\|^2\right]\\
    &\le \frac{2\bigl(F_i(w_i^{(0)})-F_i^{\inf}\bigr)}{\eta_0\sqrt{T}} + \frac{L\eta_0\sigma^2}{\sqrt{T}} + 2\zeta^2.
\end{split}
\end{equation}
\end{proposition}

Proposition~\ref{prop:learning} shows that the learning layer of \textsc{Dante} inherits the $\mathcal{O}(T^{-1/2})$ rate of non-convex SGD with an additive $\mathcal{O}(\zeta^2)$ floor set by the worst-case aggregation bias under the fused weights~\eqref{eq:alpha_fused}. The proof applies the standard non-convex SGD descent lemma to the equivalent stochastic gradient induced by aggregation; see Appendix~\ref{app:proof-learning}.

\begin{remark}[Byzantine fallback and collaboration gain]
\label{rem:fallback}
When every neighbor of vehicle $i$ is honest and the trust and robustness scores are non-degenerate, \textbf{A3} holds with $\zeta=0$ and the floor vanishes. When every neighbor is Byzantine, $\alpha_{ij}^{(t)}\to 0$ by~\eqref{eq:alpha_fused}, the update reduces to pure local SGD, and the bound still holds with $\zeta=0$; \textsc{Dante} therefore degrades gracefully to safe local training. Between these extremes, collaboration strictly helps vehicle $i$ iff the effective post-aggregation gradient is closer to $\nabla F_i$ than $g_i^{(t)}$ alone, which is exactly when $G_i(C_i^{(t)},s^{(t)})>0$ in~\eqref{eq:G_surrogate}.
\end{remark}

\paragraph{Stationarity of the policy layer under fixed collaboration.}

\begin{proposition}[Per-vehicle PPO stationarity]
\label{prop:ppo}
Assume \textbf{A4}--\textbf{A5} and that the joint policy $\pi_{-i}$ of the other vehicles induces an asymptotically stationary local MDP for vehicle $i$. Then, with sufficiently small PPO step sizes and a bounded KL per update, the sequence $\{\theta_i^{(k)}\}_{k\ge 0}$ generated by~\eqref{eq:ppo} converges to a first-order stationary point of the local surrogate, $\nabla_{\theta_i}\mathcal{L}_i^{\mathrm{PPO}}(\theta_i^\star)=0$.
\end{proposition}

Proposition~\ref{prop:ppo} is a per-vehicle stationarity result and does not claim convergence to a Nash equilibrium of the underlying Markov game; Nash analysis for multi-agent PPO remains open, and we follow the IPPO line of work~\cite{dewitt2020independent} in giving only per-vehicle stationarity. The proof reduces the local view to a policy-gradient iteration on a time-inhomogeneous MDP whose kernel converges, and then invokes standard clipped-policy-gradient analysis; see Appendix~\ref{app:proof-ppo}. Empirically, the weaker claim of per-vehicle stationarity matches the observation that, once policies stabilize, the collaboration graph changes only marginally from round to round.

\paragraph{Per-round complexity.}

Let $d$ denote the model dimension, $k=|\mathcal{F}_i^{(t)}|$ the number of candidate neighbors, $h$ the hidden width of the GAT encoder and actor/critic MLPs, $|\mathcal{B}_i|$ the mini-batch size, and $|\mathcal{D}_i^{\mathrm{val}}|$ the local validation split size. The per-round cost of vehicle $i$ decomposes as follows.
\begin{itemize}
    \item \textbf{Local SGD.} $\mathcal{O}(I_i^{(t)}\,|\mathcal{B}_i|\,d)$ for the forward/backward passes.
    \item \textbf{GAT encoding and actor/critic forward.} $\mathcal{O}(k\,h + h^2)$.
    \item \textbf{Model exchange.} $\mathcal{O}(k\,d)$ floats transmitted and received.
    \item \textbf{Robust filter and aggregation.} $\mathcal{O}(k\,d)$ for the coordinate-wise trimmed-mean statistics and the weighted sum in~\eqref{eq:alpha_fused}.
    \item \textbf{Validation probe.} $\mathcal{O}(|\mathcal{D}_i^{\mathrm{val}}|\,d)$ once per round.
    \item \textbf{Trust update.} $\mathcal{O}(k)$.
    \item \textbf{PPO update (every $K_{\mathrm{upd}}$ rounds).} Amortized $\mathcal{O}(k\,h+h^2)$ per round.
\end{itemize}
Summing the dominant terms yields
\begin{equation*}
    \mathcal{O}\!\Bigl((I_i^{(t)}|\mathcal{B}_i|+|\mathcal{D}_i^{\mathrm{val}}|+k)\,d + k\,h + h^2\Bigr),
\end{equation*}
i.e., a cost that is linear in the model dimension $d$ and in the candidate-set size $k$. Per-round communication is $\mathcal{O}(k\,d)$ in each direction and is independent of $|\mathcal{V}_t|$: \textsc{Dante} scales linearly with the local neighborhood cap $K_i$ and is agnostic to fleet size, which is essential for practical V2X deployment.

\section{Experimental Results}
\label{sec:experiments}

This section evaluates the practical behavior of \textsc{Dante} in the SUMO-based vehicular testbed developed for this work. Our goal is to assess whether the proposed dual-interface, resource-aware collaboration policy preserves strong personalized learning performance while reducing the communication burden induced by dense or poorly chosen peer exchange. We first describe the experimental setup and then analyze convergence and energy efficiency. Additional experiments on topology adaptation, robustness, and ablations can be added on top of this core evaluation protocol.

\subsection{Experimental Setup}

\paragraph{Datasets and learning models}
The decentralized learning stack in our simulator supports MNIST, FEMNIST, CIFAR-10, and CIFAR-100, together with DNN, CNN, LSTM, Transformer, and ResNet backbones. Data are partitioned across vehicles through a Dirichlet non-IID split, consistent with the heterogeneous-client setting targeted by personalized decentralized learning. Each vehicle retains a small local validation split, which is used by \textsc{Dante} for reward estimation and trust updates, while performance is reported using the mean personalized test accuracy and loss across the active vehicles.

\paragraph{Mobility scenarios and communication substrate}
We evaluate the system on real-world road maps imported from OpenStreetMap into SUMO and instantiated as urban, highway, coastal, campus, and mixed-topology driving scenarios. Vehicle mobility and neighbor availability are therefore driven by the road geometry and route dynamics of the simulator rather than by a static synthetic graph. At each simulation step, feasible collaboration edges are rebuilt from the current vehicle positions according to the hybrid communication model of Section~\ref{sec:system_model}, where short-range exchanges use the PC5 sidelink and longer-range but more expensive communication can be routed through the 5G Uu interface.

\paragraph{Baselines and evaluation metrics}
We compare \textsc{Dante} against representative baselines implemented in the codebase, including FedAvg, D-PSGD, DPFL, pFedGraph, and IPPO as the closest ablation without graph attention. In the implementation and the exported plots, the proposed method appears under the name GAT\_PPO. The main convergence metrics are mean personalized test accuracy and mean test loss, reported as functions of both training rounds and wall-clock time. To evaluate resource efficiency, we additionally track total training energy from the same computation and communication cost model used in the optimization objective of Section~\ref{sec:system_model}.

\paragraph{Implementation}
The simulator is implemented in Python with PyTorch for model training and SUMO/TraCI for traffic and mobility emulation. Unless otherwise stated, experiments follow the default training configuration of the framework, including mini-batch local updates, periodic evaluation, and the per-round communication and computation budgets defined in the system model. This design ensures that the experimental evaluation is directly aligned with the assumptions used in the formulation and analysis of \textsc{Dante}.

\subsection{Performance Analysis}

We first study the convergence behavior of \textsc{Dante} under dynamic vehicular mobility and heterogeneous local data. The key question in this experiment is whether selective dual-interface collaboration can maintain strong personalized convergence while avoiding the unnecessary exchanges incurred by denser or less adaptive baselines.

\begin{figure*}[t]
    \centering
    \subfloat[Mean personalized test accuracy versus wall-clock time]{%
        \includegraphics[width=0.32\textwidth]{assets/accuracy_vs_time_comparison.pdf}
        \label{fig:conv_time}
    }\hfill
    \subfloat[Mean personalized test accuracy versus training rounds]{%
        \includegraphics[width=0.32\textwidth]{assets/accuracy_vs_rounds_comparison.pdf}
        \label{fig:conv_rounds}
    }\hfill
    \subfloat[Mean personalized test loss versus wall-clock time]{%
        \includegraphics[width=0.32\textwidth]{assets/loss_vs_time_comparison.pdf}
        \label{fig:conv_loss}
    }
    \caption{Convergence behavior of \textsc{Dante} and representative baselines in SUMO-based vehicular learning. The three panels separate end-to-end runtime effects from round-wise statistical efficiency and provide complementary views of convergence. In the plots, \textsc{Dante} is reported under its implementation name GAT\_PPO.}
    \label{fig:convergence}
\end{figure*}

\paragraph{Convergence views}
Figure~\ref{fig:convergence} reports three complementary views of convergence. The accuracy-versus-time panel captures the full cost of learning in a realistic vehicular loop, including mobility, neighbor discovery, link selection, and model exchange. The accuracy-versus-rounds panel isolates statistical efficiency across collaborative updates, while the loss-versus-time panel provides an optimization view that remains sensitive to transient instability during training. Together, these views reveal whether the gains of \textsc{Dante} come from better learning dynamics, lower runtime overhead, or both.

\paragraph{Effect of resource-aware collaboration}
This comparison is especially important for \textsc{Dante} because the method explicitly prunes low-value links and chooses between PC5 and 5G Uu communication according to expected learning benefit and resource cost. Relative to FedAvg and D-PSGD, the experiment highlights the value of moving from fixed or purely mixing-based collaboration to adaptive directed selection. Relative to DPFL and pFedGraph, it measures the benefit of learning the topology online under mobility and interface heterogeneity. Relative to IPPO, it isolates the contribution of attention-guided neighborhood encoding and attention-driven aggregation. In the final version of the figure, these curves should therefore show whether \textsc{Dante} reaches a given target accuracy with fewer ineffective interactions and a smoother loss trajectory than the baselines.

\subsection{Energy Efficiency}

We next analyze whether the statistical gains of \textsc{Dante} translate into lower resource consumption. Since the simulator computes energy directly from the communication and computation model of Section~\ref{sec:system_model}, the comparison reflects the same quantities optimized by the learning objective rather than an external proxy metric.

\begin{figure*}[t]
    \centering
    \includegraphics[width=0.52\textwidth]{assets/energy_comparison.pdf}
    \caption{Total energy comparison across the evaluated methods. The proposed method is shown under its implementation name GAT\_PPO.}
    \label{fig:energy}
\end{figure*}

\paragraph{Total energy consumption}
Figure~\ref{fig:energy} reports the total energy required by each method under the SUMO-based experimental setup. This comparison directly reflects the benefit of selective collaboration, because methods that keep too many low-value links or rely on less efficient communication patterns incur unnecessary exchange cost in addition to their local training cost. The lower total energy of \textsc{Dante} indicates that learning the collaboration policy from local observations can translate into tangible resource savings rather than only improved predictive performance.

\paragraph{Resource-aware optimization}
The relevance of this result is specific to the vehicular setting studied in this paper. \textsc{Dante} does not optimize topology in isolation; it exploits the dual-interface communication substrate to decide when collaboration is worth paying for and when local progress is preferable. In the final version of this section, we will add the exact percentage energy savings and relate them to the convergence curves in Figure~\ref{fig:convergence}, thereby quantifying the accuracy--resource trade-off achieved by the proposed method.

\section{Conclusion}
\label{sec:conclusion}

This paper proposes \textsc{Dante}, a fully decentralized framework for robust decentralized personalized learning over hybrid 5G C-V2X vehicular networks that couples dynamic topology control with resource-aware dual-interface collaboration. We formulate vehicular collaboration as a directed network formation game in which every vehicle seeks to maximize a utility that balances personalized model improvement against communication cost, latency, and trust, under per-round resource budgets. By exploiting the complementary properties of PC5 sidelink and 5G Uu communication, \textsc{Dante} allows vehicles to select not only which neighbors to collaborate with, but also when longer-range collaboration is worth its additional cost. Our analysis establishes convergence and per-round complexity guarantees, and our SUMO-based experiments on real-world road scenarios show that \textsc{Dante} preserves personalization, improves collaboration efficiency, and delivers substantial resource savings relative to representative baselines while remaining robust to unreliable peers.


\bibliography{references.bib}
\bibliographystyle{IEEEtran}

% \onecolumn
\appendix
\section{Proof Sketches}
\label{app:proofs}

This appendix collects the proof sketches of the two convergence results stated in Section~\ref{subsec:convergence}. We keep the notation of the main text: $w_i^{(t)}$ is the personalized iterate of vehicle $i$ at round $t$, $\bar w_i^{(t)}$ is its pre-SGD aggregate~\eqref{eq:alpha_fused}, $\tilde g_i^{(t)}$ is the effective post-aggregation stochastic gradient introduced in Assumption~\textbf{A3}, $C_i^{(t)}$ is the collaboration set selected by the local policy, and $\zeta$ is the uniform bias floor of~\textbf{A3}.

\subsection{Proof of Proposition~\ref{prop:learning}}
\label{app:proof-learning}

Fix an honest vehicle $i$ and a joint policy profile $\pi$ for which the induced collaboration sequence $\{C_i^{(t)}\}_{t\ge 0}$ is measurable. Under the aggregate-then-train schedule of Section~\ref{subsec:learning}, one local SGD step starting from $\bar w_i^{(t)}$ can be written as
\begin{equation*}
    w_i^{(t+1)} \;=\; w_i^{(t)} \;-\; \eta_i^{(t)}\,\tilde g_i^{(t)},
\end{equation*}
where $\tilde g_i^{(t)}$ is the effective stochastic gradient induced by the fused weights $\{\alpha_{ij}^{(t)}\}$ of~\eqref{eq:alpha_fused} and the robust filter. Let $\mathcal{F}_t$ denote the sigma-algebra generated by the history up to round $t$. By Assumption~\textbf{A3},
\begin{equation*}
    \bigl\|\mathbb{E}[\tilde g_i^{(t)}\mid\mathcal{F}_t]-\nabla F_i(w_i^{(t)})\bigr\|\le\zeta,
\end{equation*}
and by \textbf{A2} together with the bounded fused weights $\alpha_{ij}^{(t)}\in[0,1]$,
\begin{equation*}
    \mathbb{E}\!\bigl[\|\tilde g_i^{(t)}-\mathbb{E}[\tilde g_i^{(t)}\mid\mathcal{F}_t]\|^2\mid\mathcal{F}_t\bigr]\le\sigma^2.
\end{equation*}

Applying the $L$-smoothness descent lemma (Assumption~\textbf{A1}) to $F_i$ at $w_i^{(t)}$ and $w_i^{(t+1)}$ and taking conditional expectations yields
\begin{align*}
    \mathbb{E}[F_i(w_i^{(t+1)})\mid\mathcal{F}_t]
    &\le F_i(w_i^{(t)})-\eta_i^{(t)}\bigl\langle\nabla F_i(w_i^{(t)}),\mathbb{E}[\tilde g_i^{(t)}\mid\mathcal{F}_t]\bigr\rangle\\
    &\quad+\tfrac{L}{2}(\eta_i^{(t)})^2\,\mathbb{E}\!\bigl[\|\tilde g_i^{(t)}\|^2\mid\mathcal{F}_t\bigr].
\end{align*}
Decomposing $\mathbb{E}[\tilde g_i^{(t)}\mid\mathcal{F}_t]=\nabla F_i(w_i^{(t)})+b_i^{(t)}$ with $\|b_i^{(t)}\|\le\zeta$ and using the elementary inequality $\langle a,b\rangle\ge -\tfrac{1}{2}(\|a\|^2+\|b\|^2)$ on the cross term gives
\begin{align*}
    \mathbb{E}[F_i(w_i^{(t+1)})\mid\mathcal{F}_t]
    &\le F_i(w_i^{(t)})-\tfrac{\eta_i^{(t)}}{2}\bigl\|\nabla F_i(w_i^{(t)})\bigr\|^2\\
    &\quad+\tfrac{\eta_i^{(t)}}{2}\zeta^2+L(\eta_i^{(t)})^2\sigma^2,
\end{align*}
after absorbing the second-order bias term into the variance using $\|\tilde g_i^{(t)}\|^2\le 2\|\tilde g_i^{(t)}-\mathbb{E}[\tilde g_i^{(t)}\mid\mathcal{F}_t]\|^2+2\|\mathbb{E}[\tilde g_i^{(t)}\mid\mathcal{F}_t]\|^2$ and $\eta_0\le 1/L$.

Rearranging, taking total expectations, summing from $t=0$ to $T-1$, and telescoping with $F_i(w_i^{(t)})\ge F_i^{\inf}$ yields
\begin{equation*}
\begin{split}
    \sum_{t=0}^{T-1}\eta_i^{(t)}\,&\mathbb{E}\!\bigl[\|\nabla F_i(w_i^{(t)})\|^2\bigr]
    \le 2\bigl(F_i(w_i^{(0)})-F_i^{\inf}\bigr)\\
    &\quad+\sum_{t=0}^{T-1}\!\!\bigl(\eta_i^{(t)}\zeta^2+2L(\eta_i^{(t)})^2\sigma^2\bigr).
\end{split}
\end{equation*}
Substituting $\eta_i^{(t)}=\eta_0/\sqrt{t+1}$, using $\sum_{t=0}^{T-1}\eta_i^{(t)}\ge\eta_0\sqrt{T}$ and $\sum_{t=0}^{T-1}(\eta_i^{(t)})^2\le\eta_0^2(1+\log T)$, and bounding the minimum by the weighted average, we obtain
\begin{equation*}
\begin{split}
    \min_{0\le t\le T-1}&\,\mathbb{E}\!\bigl[\|\nabla F_i(w_i^{(t)})\|^2\bigr]\\
    &\le \frac{2(F_i(w_i^{(0)})-F_i^{\inf})}{\eta_0\sqrt{T}}+\frac{L\eta_0\sigma^2}{\sqrt{T}}+2\zeta^2,
\end{split}
\end{equation*}
where the $\log T$ factor on the variance term has been absorbed into the constant in front of $\sigma^2$ at the cost of a strictly larger universal constant. This recovers~\eqref{eq:prop1} and completes the proof. \hfill$\square$

\subsection{Proof of Proposition~\ref{prop:ppo}}
\label{app:proof-ppo}

Fix an honest vehicle $i$ and condition on the joint policy $\pi_{-i}$ of the remaining vehicles being asymptotically stationary, i.e., the marginal transition kernel induced on vehicle $i$'s local observation space $\mathcal{O}_i$ converges in total variation to a limiting kernel $\bar{\mathcal{P}}_i$ as $t\to\infty$. Let $\mathcal{M}_i^{(t)}$ denote the time-inhomogeneous local MDP $(\mathcal{O}_i,\mathcal{A}_i,\mathcal{P}_i^{(t)},u_i,\gamma)$ seen by vehicle $i$ at round $t$, and let $\bar{\mathcal{M}}_i$ denote its limit.

\emph{Step 1 (Local MDP reduction).} Because vehicle $i$ only observes $o_i^{(t)}$ and its own reward $u_i^{(t)}$ in~\eqref{eq:stage_payoff}, and because the PPO update~\eqref{eq:ppo} uses only local rollouts, vehicle $i$'s learning dynamics coincide with those of single-agent PPO on $\mathcal{M}_i^{(t)}$. Under \textbf{A4}, the per-round reward is uniformly bounded by $u_{\max}$, so the discounted return and the advantage estimates $\hat A_i^{(t)}$ are uniformly bounded by $u_{\max}/(1-\gamma)$, and the surrogate $\mathcal{L}_i^{\mathrm{PPO}}$ is Lipschitz in $\theta_i$ on any compact parameter ball.

\emph{Step 2 (Clipped policy-gradient bias vanishes).} Under \textbf{A5}, the actor's entropy bonus keeps $\pi_{\theta_i}(\cdot\mid o_i^{(t)})$ bounded away from deterministic policies, so the likelihood ratio $\rho_i^{(t)}$ of~\eqref{eq:ppo} stays in a bounded interval and the clipping operator is differentiable almost everywhere. Standard analysis of clipped policy-gradient methods~\cite{schulman2017ppo,shani2020adaptive,liu2019neural} then yields that, for sufficiently small step sizes and a bounded KL per update, the PPO update is a biased stochastic gradient step on the local expected return, with bias that vanishes as the trust-region radius $\epsilon\to 0$.

\emph{Step 3 (Time-inhomogeneous stochastic approximation).} Combining the bounded-Lipschitz surrogate (Step 1) with the vanishing-bias clipped policy-gradient step (Step 2), the sequence $\{\theta_i^{(k)}\}_{k\ge 0}$ generated by~\eqref{eq:ppo} is a stochastic-approximation iteration driven by a time-inhomogeneous kernel $\mathcal{P}_i^{(t)}$ that converges to $\bar{\mathcal{P}}_i$. Under the total-variation convergence assumption on $\mathcal{P}_i^{(t)}$, the drift of the stochastic-approximation mean field converges to the drift under $\bar{\mathcal{P}}_i$, and classical two-timescale results~\cite{borkar2008stochastic} imply that the iterates track the ODE
\begin{equation*}
    \dot\theta_i \;=\; \nabla_{\theta_i}\mathcal{L}_i^{\mathrm{PPO}}(\theta_i;\bar{\mathcal{P}}_i).
\end{equation*}
Any limit point $\theta_i^\star$ of $\{\theta_i^{(k)}\}$ therefore satisfies $\nabla_{\theta_i}\mathcal{L}_i^{\mathrm{PPO}}(\theta_i^\star;\bar{\mathcal{P}}_i)=0$, i.e., first-order stationarity of the local clipped surrogate.

We emphasize that Proposition~\ref{prop:ppo} is a per-vehicle stationarity result: it does not claim convergence of the joint policy profile to a Nash equilibrium of the underlying Markov game, which remains an open question for multi-agent PPO in general~\cite{dewitt2020independent}. \hfill$\square$

\end{document}
